% =====================================================================
%  Chương 7 — Deep Graph và mạng nơ-ron đồ thị
% =====================================================================
\chapter{Deep Graph và mạng nơ-ron đồ thị}\label{chap:gnn}

Các chương trước làm việc trên bảng, lưới ảnh và chuỗi, nơi mỗi mẫu có vị trí cố định. Bài báo trích dẫn nhau, sản phẩm được
mua cùng nhau hay tiền chạy qua các giao dịch thì không có thứ tự như vậy: số láng giềng thay đổi tùy nút, và đánh số lại nút
không làm đồ thị khác đi. Chương này xây GCN, GraphSAGE, GAT trên ba đồ thị thật, đo chúng với hai đường cơ sở, rồi dự đoán
liên kết trên đồ thị Hetionet đã nạp vào Neo4j ở Chương~5.

\section{Cơ sở toán học}

Đồ thị vô hướng $G=(V,E)$ có $N$ nút được mô tả bởi ma trận kề $A$ ($A_{ij}=1$ khi có cạnh), ma trận bậc đường chéo
$D_{ii}=\sum_j A_{ij}$ và ma trận đặc trưng $X\in\mathbb R^{N\times F}$ (Hình~\ref{fig:ch7-repr}). Laplacian $L=D-A$ và dạng
chuẩn hóa đối xứng $L_{\mathrm{sym}}=I-D^{-1/2}AD^{-1/2}$ đo độ gồ ghề của tín hiệu $x$ trên đồ thị:
\begin{equation}
x^{\top}L_{\mathrm{sym}}\,x=\sum_{(i,j)\in E}\Big(\frac{x_i}{\sqrt{D_{ii}}}-\frac{x_j}{\sqrt{D_{jj}}}\Big)^2 .
\end{equation}
$L_{\mathrm{sym}}$ đối xứng nửa xác định dương, trị riêng nằm trong $[0,2]$, số trị riêng 0 bằng số thành phần liên thông.
Lọc tín hiệu theo phổ cần phân rã trị riêng $O(N^3)$; Kipf và Welling xấp xỉ bậc một rồi thêm tự khuyên
$\tilde A=A+I$, $\tilde D=D+I$, được phép truyền một lớp chỉ tốn $O(|E|)$~\cite{ch7:kipf2017}:
\begin{equation}
H^{(\ell+1)}=\sigma\big(\hat A\,H^{(\ell)}W^{(\ell)}\big),\qquad \hat A=\tilde D^{-1/2}\tilde A\tilde D^{-1/2}.
\end{equation}
Chuẩn hóa hai phía làm trọng số cạnh bằng $1/\sqrt{\tilde d_i\tilde d_j}$: nút trung tâm nhiều cạnh không áp đảo láng giềng,
còn tự khuyên giữ lại thông tin của chính nút.

\hinhtikz{ch7_graph_repr}{Bốn cách nhìn một đồ thị 5 nút: cấu trúc, ma trận kề, ma trận bậc và Laplacian chuẩn hóa.}{ch7-repr}

\textbf{Truyền thông điệp (MPNN).} Gilmer và cộng sự gộp phần lớn GNN vào một khuôn~\cite{ch7:gilmer2017}
(Hình~\ref{fig:ch7-mpnn}): mỗi cạnh sinh thông điệp $\psi$, mỗi nút gộp thông điệp bằng phép $\bigoplus$ không phụ thuộc thứ tự
(tổng, trung bình, cực đại), rồi cập nhật trạng thái bằng $\phi$. Sau $L$ lớp, một nút nghe được vùng cách nó $L$ bước.
Ba kiến trúc khác nhau ở cách gán trọng số cho láng giềng:
\begin{align}
\text{GCN:}\quad & h_v'=\sigma\Big(\sum_{u\in\mathcal N(v)\cup\{v\}}\frac{1}{\sqrt{\tilde d_u\tilde d_v}}\,Wh_u\Big),\\
\text{GraphSAGE:}\quad & h_v'=\sigma\big(W_1h_v+W_2\,\mathrm{mean}\{h_u: u\in\mathcal S(v)\}\big),\\
\text{GAT:}\quad & h_v'=\big\Vert_{k=1}^{K}\sigma\Big(\sum_{u}\alpha^{k}_{vu}W^{k}h_u\Big).
\end{align}
GCN cố định trọng số theo bậc. GraphSAGE~\cite{ch7:hamilton2017} tách riêng trọng số của chính nút và lấy mẫu tập láng giềng
$\mathcal S(v)$ lúc học, nên áp được cho nút chưa từng thấy. GAT~\cite{ch7:velickovic2018} học hệ số chú ý $\alpha_{vu}$ từ
chính đặc trưng hai đầu cạnh (Hình~\ref{fig:ch7-gat}).

\hinhtikz{ch7_mpnn3d}{Truyền thông điệp qua ba lớp: cùng một đồ thị, nút đỏ gộp thông tin láng giềng ở lớp trước.}{ch7-mpnn}

\hinhtikz{ch7_gcn3d}{GCN hai lớp cho phân loại nút: $Z=\hat A\,\mathrm{ReLU}(\hat AXW^{(0)})W^{(1)}$.}{ch7-gcn}

\hinhtikz{ch7_gat}{Tính hệ số chú ý của GAT và ghép nhiều đầu chú ý.}{ch7-gat}

\section{Ba đồ thị thực nghiệm}

\textbf{Nguồn.} \emph{Cora} lấy từ kho Planetoid~\cite{ch7:yang2016}
(\url{https://github.com/kimiyoung/planetoid}, bản dữ liệu dùng lại từ Tiểu luận 03): nút là bài báo, cạnh là trích dẫn,
đặc trưng là túi từ nhị phân. \emph{Amazon Photo} của Shchur và cộng sự~\cite{ch7:shchur2018}
(\url{https://github.com/shchur/gnn-benchmark}): nút là sản phẩm, cạnh nối hai sản phẩm hay được mua cùng nhau; bộ gốc không
công bố tên lớp nên ghi Lớp 0 đến Lớp 7. \emph{Elliptic Bitcoin}~\cite{ch7:weber2019}, tải bản sao của PyG
(\url{https://data.pyg.org/datasets/elliptic/}): nút là giao dịch, cạnh là dòng tiền, đặc trưng là 165 thuộc tính đã ẩn danh.
Tôi chọn Elliptic thay ogbn-arxiv vì bộ này có chiều thời gian thật và lệch lớp nặng, tức hai cái bẫy đánh giá mà Cora
không có. Cả ba dùng toàn bộ đồ thị. Bảng~\ref{tab:ch7-data} là số đo thật trên dữ liệu đã tải.

\bang{ch7_datasets}{Kích cỡ và tính chất đồ thị. Bậc, đồng dạng tính sau khi đối xứng hóa cạnh; đồng dạng trên Elliptic chỉ tính cạnh có nhãn ở hai đầu.}{ch7-data}

\hinhpdf{ch7_data}{1.0}{Phân bố lớp (hàng trên) và phân bố bậc thang log-log (hàng dưới). Elliptic chỉ đếm nút có nhãn.}{ch7-dist}

\bang{ch7_samples}{Nút mẫu thật của Cora và Amazon Photo: nút đầu tiên của hai lớp đầu.}{ch7-mau}

\bang{ch7_samples_ell}{Giao dịch mẫu thật của Elliptic. Đặc trưng đã z-score bằng thống kê các bước 1 đến 29.}{ch7-mau-ell}

\nhanxet{Đồng dạng cạnh (homophily) của Cora và Photo lần lượt là \SL{ch7_data/cora/edge_homophily} và
\SL{ch7_data/photo/edge_homophily}, nên trung bình láng giềng là phép lọc có lợi. Đuôi bậc dài: Cora có nút bậc
\SL{ch7_data/cora/deg_max}, Photo có nút bậc \SL{ch7_data/photo/deg_max} cùng \SL{ch7_data/photo/isolated} nút cô lập mà GNN
chỉ thấy đặc trưng của chính chúng. Elliptic có đồng dạng thô \SL{ch7_data/elliptic/edge_homophily} nhưng chỉ
\SL{ch7_data/elliptic/adj_homophily} sau hiệu chỉnh, vì \SL{ch7_data/elliptic/imbalance} giao dịch licit mới có một illicit nên
hai đầu cạnh trùng nhãn phần nhiều do tình cờ. Chỉ \SL{ch7_data/elliptic/labeled} trên \SL{ch7_data/elliptic/nodes} nút có nhãn.
Đồ thị này tách thành đúng \SL{ch7_data/elliptic/components} thành phần liên thông ứng với \SL{ch7_data/elliptic_time/steps}
bước thời gian, số cạnh vắt qua hai bước là \SL{ch7_data/elliptic_time/cross_step_edges}.}

\textbf{Tiền xử lý.} Cạnh được đối xứng hóa, bỏ tự khuyên. Cora, Photo chia mỗi hàng đặc trưng cho tổng của hàng, phép biến
đổi theo từng nút nên không học gì từ test. Elliptic z-score bằng trung bình, độ lệch của riêng các bước 1 đến 29.
Cora dùng phép chia chuẩn \SL{ch7_data/cora/train}/\SL{ch7_data/cora/val}/\SL{ch7_data/cora/test}. Photo không có phép chia
chuẩn nên lấy 20 nút mỗi lớp cho train, 30 cho val, còn lại test, mỗi seed một phép chia. Elliptic chia theo thời gian:
\SL{ch7_data/elliptic/train} nút có nhãn ở bước 1--29, \SL{ch7_data/elliptic/val} ở bước 30--34, \SL{ch7_data/elliptic/test}
ở bước 35--49, kèm trọng số lớp nghịch đảo tần suất trong hàm mất mát.

\section{Hiện thực: GCN từ đầu và PyTorch thuần}

Bản NumPy không dùng autograd hay ma trận thưa: $\hat A$ là mảng dày $N\times N$, đạo hàm suy tay theo quy tắc chuỗi. Với
$G=\partial\mathcal L/\partial Z$ ta có $\partial\mathcal L/\partial W^{(1)}=H^{(1)\top}\hat A^{\top}G$, rồi đi ngược qua dropout
và ReLU để được $\partial\mathcal L/\partial W^{(0)}$ (Đoạn mã~\ref{code:ch7_numpy__GCNNumpy.loss_grads}). Hàm \texttt{grad\_check}
so đạo hàm này với sai phân trung tâm trên đồ thị ngẫu nhiên 12 nút, sai lệch lớn nhất \SL{ch7_numpy/grad_check_max_error}.

\maNguon{ch7_numpy__GCNNumpy.loss_grads}{Lan truyền ngược viết tay của GCN NumPy.}

\maNguon{ch7_gnn__gcn_norm_adj}{Tính $\hat A$ thưa bằng PyTorch thuần, dùng cho mọi lớp GCN.}

\maNguon{ch7_gnn__GATLayer.forward}{Lớp GAT: điểm cạnh, softmax theo nút đích, gộp có trọng số.}

\begin{ketqua}{Đầu ra thực của GCN NumPy trên Cora: ba seed đầu và trung bình \SL{ch7_numpy/n_seeds} seed}
numpy 0 0.817 26.7 s\\
numpy 1 0.822 25.7 s\\
numpy 2 0.814 25.8 s\\
trung bình độ đúng \SL{ch7_numpy/accuracy}, F1 macro \SL{ch7_numpy/f1_macro}
\end{ketqua}

\section{Đánh giá ba mô hình trên ba đồ thị}

\subsection{So sánh với đường cơ sở}

MLP chỉ dùng đặc trưng; Label Propagation (LP) chỉ dùng cạnh, lặp $F\leftarrow\alpha SF+(1-\alpha)Y$~\cite{ch7:zhou2004} với $\alpha$
chọn trên val. Mô hình hai lớp, ẩn 64 chiều (Elliptic 128), Adam, chọn trọng số theo F1 macro của val, dừng sớm; riêng test chỉ đo
một lần ở cuối. GraphSAGE lấy mẫu tối đa 10 láng giềng (Photo 25) khi học. Các siêu tham số Photo, Elliptic được chỉnh trên val.
Mỗi cấu hình chạy \SL{ch7_node/cora/gcn/n_seeds} seed trên GPU.

\bang{ch7_compare}{Phân loại nút, trung bình trên các seed. Elliptic đo bằng chia thời gian, cột F1 illicit là thước đo thật của bài toán.}{ch7-compare}

\hinhpdf{ch7_compare}{1.0}{F1 macro của năm phương pháp trên ba đồ thị, thanh sai số là độ lệch chuẩn giữa các seed.}{ch7-compare}

\begin{phathien}{Đọc kết quả}
Trên Cora, GCN đạt F1 macro \SL{ch7_node/cora/gcn/f1_macro}, hơn MLP \SL{ch7_summary/cora/gcn_vs_mlp}; bản NumPy đạt
\SL{ch7_numpy/f1_macro}, khớp bản PyTorch, cho thấy đạo hàm viết tay đúng. Trên Photo, \SL{ch7_summary/photo/best} dẫn đầu với
\SL{ch7_summary/photo/best_f1}, hơn MLP \SL{ch7_summary/photo/gain_vs_mlp}. LP không dùng đặc trưng mà vẫn đạt
\SL{ch7_node/cora/lp/f1_macro} trên Cora và \SL{ch7_node/photo/lp/f1_macro} trên Photo, nghĩa là riêng cấu trúc đã mang nhiều
thông tin nhãn. Trên Elliptic, LP cho F1 illicit bằng \SL{ch7_node/elliptic/lp/illicit_f1} vì nhãn lan tới \SL{ch7_node/elliptic/lp/reached}
phần nút test: không cạnh nào nối hai bước thời gian, nên nhãn quá khứ không có đường đi tới tương lai, và độ đúng
\SL{ch7_node/elliptic/lp/accuracy} của nó chỉ là tỷ lệ lớp đa số.
\end{phathien}

Trên Elliptic, mọi GNN đều học quy nạp vì đồ thị test không dính với đồ thị train. Tốt nhất là \SL{ch7_summary/elliptic/best}, F1
illicit \SL{ch7_node/elliptic/sage/illicit_f1}; GAT chỉ đạt \SL{ch7_node/elliptic/gat/illicit_f1}, thấp hơn cả MLP
\SL{ch7_node/elliptic/mlp/illicit_f1} dù tốn thời gian nhất. Hình~\ref{fig:ch7-elliptic} cho biết lý do tổng thể không cao: F1
trung bình của GCN là \SL{ch7_summary/elliptic/before43/gcn} ở các bước 35--42 rồi rơi xuống
\SL{ch7_summary/elliptic/after43/gcn} từ bước 43. Weber và cộng sự cũng ghi nhận cú rơi tại đây, trùng thời điểm một chợ đen lớn
đóng cửa~\cite{ch7:weber2019}: kiểu gian lận đổi, và không mô hình nào trong năm mô hình theo kịp.

\hinhpdf{ch7_elliptic}{1.0}{Trái: số giao dịch illicit theo bước, nét đứt ngăn train, val, test. Phải: F1 illicit trên test theo bước, trung bình các seed.}{ch7-elliptic}

\hinhpdf{ch7_err}{1.0}{Phân tích lỗi: độ đúng test theo tỷ lệ láng giềng cùng nhãn của nút (seed 0).}{ch7-err}

\nhanxet{Phân tích lỗi cho thấy GNN giỏi ở đâu và hỏng ở đâu (Hình~\ref{fig:ch7-err}). Ở nhóm nút có trên 75\,\% láng giềng cùng nhãn,
GCN đạt \SL{ch7_summary/cora/high_gcn} trên Cora, so với MLP \SL{ch7_summary/cora/high_mlp}. Ở nhóm dưới 25\,\%, tình thế đảo
ngược: GCN \SL{ch7_summary/cora/low_gcn} so với MLP \SL{ch7_summary/cora/low_mlp} trên \SL{ch7_summary/cora/low_n} nút Cora, và
\SL{ch7_summary/photo/low_gcn} so với \SL{ch7_summary/photo/low_mlp} trên \SL{ch7_summary/photo/low_n} nút Photo. Láng giềng khác
nhãn kéo biểu diễn về lớp sai, đúng cái giá phải trả của phép làm trơn.}

\subsection{Over-smoothing theo độ sâu}

Số liệu độ sâu trên Cora lấy lại nguyên từ Tiểu luận 03 (cùng dữ liệu, GCN ẩn 64, 3 seed), không chạy lại. Hai lớp đạt
\SL{ch7_summary/depth2}; tám lớp còn \SL{ch7_summary/depth8}, còn nối tắt giữ được \SL{ch7_summary/depth8_res}. Bỏ hẳn việc học, nhân
$\hat A$ lặp lại làm khoảng cách cosine trung bình MAD~\cite{ch7:chen2020} giữa các nút tụt từ \SL{ch7_summary/mad0} xuống
\SL{ch7_summary/mad64}: $\hat A^k$ hội tụ về vector riêng trội, nên mọi nút cùng thành phần liên thông dần trùng nhau
(Hình~\ref{fig:ch7-depth}).

\hinhpdf{ch7_depth}{1.0}{Over-smoothing trên Cora: độ đúng theo số lớp, và MAD của $\hat A^kX$ không qua huấn luyện.}{ch7-depth}

\subsection{Biểu diễn học được: t-SNE}

Hình~\ref{fig:ch7-tsne} chiếu đặc trưng thô và lớp ẩn GCN xuống hai chiều (3.000 nút với Photo, Elliptic; Cora lấy từ Tiểu luận 03).
Silhouette tính trên vector gốc theo nhãn thật: Cora tăng từ \SL{ch7_summary/sil/cora_raw} lên \SL{ch7_summary/sil/cora_gcn}, Photo
từ \SL{ch7_summary/sil/photo_raw} lên \SL{ch7_summary/sil/photo_gcn}, Elliptic từ \SL{ch7_summary/sil/ell_raw} lên
\SL{ch7_summary/sil/ell_gcn}.

\hinhpdf{ch7_tsne}{1.0}{t-SNE của đặc trưng thô (trên) và biểu diễn ẩn GCN (dưới), màu là nhãn thật.}{ch7-tsne}

\section{Chống rò rỉ dữ liệu trên đồ thị}

Đồ thị có ba kiểu rò rỉ mà dữ liệu bảng không có (Hình~\ref{fig:ch7-leak}). Học xuyên đồ thị cho mô hình thấy đặc trưng và cạnh của
nút test lúc truyền tin; đây là giả định hợp lệ khi đồ thị đóng, nhưng lạc quan nếu nút mới đến sau. Học quy nạp cắt mọi cạnh chạm
nút giữ lại, do \texttt{graph\_views} thực hiện (Đoạn mã~\ref{code:ch7_gnn__graph_views}): train không thấy val, test; val không thấy test.
Với dữ liệu có thời gian, chia ngẫu nhiên cho mô hình học từ tương lai.

\hinhtikz{ch7_leak}{Ba cách dựng phép chia: xuyên đồ thị, quy nạp và theo thời gian.}{ch7-leak}

\maNguon{ch7_gnn__graph_views}{Ba góc nhìn đồ thị cho học quy nạp.}

\bang{ch7_leak}{Đo rò rỉ. Cora, Photo: mức 1 là xuyên đồ thị, mức 2 là quy nạp. Elliptic: mức 1 là chia theo thời gian, mức 2 là chia ngẫu nhiên.}{ch7-leak}

\begin{phathien}{Mức lạc quan của từng phép chia}
Trên đồ thị đóng, cắt cạnh tới nút chưa thấy làm F1 macro giảm \SL{ch7_leaksum/cora_gcn_drop} với GCN trên Cora và tới \SL{ch7_leaksum/photo_sage_drop} với GraphSAGE trên Photo. Elliptic nặng hơn nhiều: chia ngẫu nhiên cho MLP F1 illicit \SL{ch7_leaksum/elliptic_mlp_random} nhưng chia theo thời gian chỉ còn \SL{ch7_leaksum/elliptic_mlp_temporal}, GCN từ \SL{ch7_leaksum/elliptic_gcn_random} còn \SL{ch7_leaksum/elliptic_gcn_temporal}. Với dữ liệu gian lận, con số từ phép chia ngẫu nhiên không nói được gì về khả năng dùng thật.
\end{phathien}

\section{Dự đoán liên kết trên Hetionet từ Neo4j}

Làn Chương 5 đã nạp tiểu đồ thị Hetionet~\cite{ch7:himmelstein2017} quanh đái tháo đường vào Neo4j, nên tôi đọc thẳng cạnh bằng Cypher
(nguồn: \SL{ch7_link/source}), không tải lại tệp. Đồ thị có \SL{ch7_link/nodes} nút thuộc \SL{ch7_link/node_kinds} kiểu,
\SL{ch7_link/triples} bộ ba thuộc \SL{ch7_link/relations} quan hệ. Bài toán: cho $(s,r,t)$, cạnh này có thật không?

\hinhtikz{ch7_linkpipe}{Quy trình dự đoán liên kết từ Neo4j tới AUC, AP.}{ch7-link}

Tách cạnh là chỗ dễ rò rỉ nhất. Hai nút có thể nối bằng nhiều quan hệ, nên tôi chia theo \emph{cặp nút} (Đoạn
mã~\ref{code:ch7_link__split_pairs}): \SL{ch7_link/train}/\SL{ch7_link/val}/\SL{ch7_link/test} bộ ba, đồ thị truyền tin chỉ gồm cạnh
train. Mẫu âm tráo đầu hoặc đuôi sang nút cùng kiểu và loại mọi bộ ba đã biết; tráo tự do sẽ dễ hơn nhiều vì chỉ cần nhận ra kiểu nút.
R-GCN~\cite{ch7:schlichtkrull2018} gán mỗi quan hệ một vector đường chéo $d_r$, giải mã bằng DistMult.

\maNguon{ch7_link__split_pairs}{Tách cạnh theo cặp nút không thứ tự.}

\bang{ch7_link}{Dự đoán liên kết, 3 seed, âm cùng kiểu. Dòng cuối cố ý để cạnh test lọt vào đồ thị truyền tin khi đo.}{ch7-linkres}

\begin{phathien}{Mô hình đơn giản thắng, và rò rỉ cho số đẹp giả}
DistMult không truyền tin đạt AUC \SL{ch7_summary/link_distmult}, cao hơn R-GCN \SL{ch7_link/clean/auc}. Với số cạnh train hiện có
nhúng nút tự do đã đủ học; hai lớp truyền tin trên đồ thị không có đặc trưng nút chưa đem lại lợi ích với số epoch
ngắn này. Heuristic Tích bậc đạt \SL{ch7_summary/link_pa}, tức một phần lớn tín hiệu chỉ là nút nào nhiều cạnh. Cùng mô hình R-GCN đã huấn
luyện, chỉ đưa cạnh test vào đồ thị truyền tin khi đo, AUC nhảy lên \SL{ch7_link/leaky/auc}, tăng \SL{ch7_summary/leak_auc_gain}: con
số đẹp nhất bảng là con số sai.
\end{phathien}

AUC theo quan hệ dao động mạnh: cao nhất là gen tham gia đường truyền (GpPW), thấp nhất là hợp chất điều trị bệnh (CtD), quan hệ có
ít cạnh test nhất. Khi xếp hạng hợp chất cho \emph{đái tháo đường type 2} bằng mô hình seed 0, trong tám ứng viên đầu ngoài cạnh
train có Sitagliptin và Miglitol, đúng là hai cạnh Hetionet đã giấu đi trong val, test; những ứng viên khác như Sildenafil
hay Drospirenone chỉ là giả thuyết của mô hình. Kết quả này chỉ mang tính học thuật, không phải khuyến nghị y khoa.

\section{Triển khai và tóm tắt}

\texttt{tlf/serve\_ch7.py} nạp đồ thị và trọng số GCN, GAT lười trên CPU, không cần PyG. Với một nút, nó trả nhãn thật, xác suất dự đoán,
đồ thị con 1 đến 2 bước, trọng số mỗi cạnh ($\alpha$ của GAT, $\hat A_{vu}$ của GCN) và danh sách thông điệp vào nút để Web App vẽ
truyền thông điệp (ảnh ở Chương~9).

Chương này xác nhận GNN vượt đường cơ sở khi cạnh đồng dạng và phép đo trung thực, đồng thời chỉ ra ba giới hạn đo được: nút lạc giữa
láng giềng khác nhãn, kiểu gian lận đổi theo thời gian, và đồ thị thưa không đặc trưng nơi mô hình nông thắng. Phép chia, không phải
kiến trúc, quyết định con số có đáng tin hay không.
